\subsection{Corte Mínimo de Aristas con Flujo Máximo (Dinic)}

\graybold{Preguntas guía:}

\begin{itemize}
\item ¿Me piden la mínima cantidad (o costo) de aristas a eliminar para separar dos nodos $s$ y $t$?
\item ¿Además de cuántas, debo imprimir CUÁLES aristas forman el corte?
\item ¿El problema se puede ver como "cuántos caminos que no comparten aristas hay de $s$ a $t$" (teorema de Menger)?
\item ¿Las aristas son no dirigidas, de modo que necesito representarlas en una red de flujo dirigida?
\end{itemize}

\graybold{Utilidad:} Calcula el flujo máximo de $s$ a $t$ en una red con capacidades, y a partir de él el corte mínimo: el conjunto de aristas de menor capacidad total cuya eliminación desconecta $t$ de $s$. Con capacidad 1 por arista, el valor es la mínima cantidad de aristas a cortar y a la vez la máxima cantidad de caminos disjuntos en aristas. El flujo es la base de emparejamientos bipartitos, asignaciones con capacidades, cortes de imágenes y problemas de selección de proyectos.

\graybold{Complejidad:} \bigO{V^2 E} en general para Dinic; con capacidades unitarias, \bigO{E \sqrt{E}}. Aquí el flujo está acotado por el grado de $s$ ($\le n - 1$) y cada camino aumentante cuesta \bigO{n + m}.

\graybold{Fundamento teórico:} El teorema de flujo máximo y corte mínimo dice que el valor del flujo máximo de $s$ a $t$ es igual a la capacidad del corte mínimo. Dinic lo calcula por fases: en cada fase, un BFS desde $s$ sobre el grafo RESIDUAL (arcos con capacidad sobrante) asigna niveles, y luego se envía flujo solo por arcos que suben exactamente un nivel, usando un puntero por nodo que avanza sobre los arcos que ya no sirven para no revisarlos de nuevo en la misma fase. Enviar flujo por un arco reduce su capacidad y aumenta la del arco reverso, lo que permite "deshacer" decisiones en caminos futuros. Guardando los arcos en arreglos de a pares, el reverso del arco $e$ es simplemente $e \oplus 1$. Una arista NO dirigida de capacidad $c$ se representa con dos arcos $a \to b$ y $b \to a$, ambos de capacidad $c$, que son el reverso uno del otro: empujar flujo en un sentido libera capacidad en el otro, lo que es exactamente el comportamiento de una arista que se puede usar en cualquier dirección. Cuando ya no hay camino de $s$ a $t$ en el residual, el flujo es máximo, y el corte se EXTRAE así: sea $S$ el conjunto de nodos alcanzables desde $s$ en el residual; toda arista original con un extremo en $S$ y el otro fuera está saturada (si no, el otro extremo sería alcanzable), y estas aristas forman un corte cuya capacidad es exactamente el flujo, así que es mínimo. En la implementación se recorren las aristas originales y se imprimen las que tienen exactamente un extremo alcanzable.

\graybold{Ejercicio aplicado}

Dado un grafo no dirigido de $n$ cruces y $m$ calles, hallar la mínima cantidad de calles a cerrar para que no haya ruta entre el cruce 1 y el cruce $n$, e imprimir cuáles.

\graybold{I/O:} $n \le 500$, $m \le 1000$ con dos enteros por calle $\Rightarrow$ lectura total + tokenización (patrón 2), separando extremos con slices de paso 2. La DFS de Dinic es recursiva: la profundidad está acotada por la longitud de un camino, a lo sumo $n = 500$, así que basta subir el límite de recursión. Como se acepta cualquier corte mínimo, la verificación comprueba que la cantidad impresa sea el mínimo real, calculado probando TODOS los subconjuntos de calles, que las calles impresas existan, y que al cerrarlas 1 y $n$ queden desconectados: se hizo en 400 grafos aleatorios, incluyendo 1 y $n$ ya desconectados, calle directa entre 1 y $n$ y varios cortes mínimos posibles, sin fallos. Con $n = 500$ y $m = 1000$ tarda \aprox{}0.01s tanto con flujo alto (1 y $n$ conectados a casi todos los nodos, corte de 498 calles) como con caminos aumentantes largos y en un grafo aleatorio.

\graybold{Código solución}

\begin{lstlisting}[language=Python]
import sys
sys.setrecursionlimit(10000)

def solve():
    data = sys.stdin.buffer.read().split()
    n = int(data[0]); m = int(data[1])
    A = list(map(int, data[2:2 + 2*m:2])); B = list(map(int, data[3:3 + 2*m:2]))
    # arcos en arreglos: el arco e y su reverso e ^ 1 van de a pares.
    # arista no dirigida = dos arcos de capacidad 1 que son reverso uno del otro
    hacia = []; cap = []
    ady = [[] for _ in range(n + 1)]
    for a, b in zip(A, B):
        ady[a].append(len(hacia)); hacia.append(b); cap.append(1)
        ady[b].append(len(hacia)); hacia.append(a); cap.append(1)
    s, t = 1, n
    flujo = 0
    while True:
        # ---- BFS: niveles en el grafo residual ----
        nivel = [-1]*(n + 1)
        nivel[s] = 0
        cola = [s]
        for u in cola:
            for e in ady[u]:
                if cap[e] and nivel[hacia[e]] < 0:
                    nivel[hacia[e]] = nivel[u] + 1
                    cola.append(hacia[e])
        if nivel[t] < 0:
            break                         # no hay camino: flujo maximo
        # ---- DFS: caminos que suben de a un nivel ----
        ptr = [0]*(n + 1)                 # arcos ya descartados en esta fase
        def dfs(u):
            if u == t:
                return True
            lista = ady[u]
            while ptr[u] < len(lista):
                e = lista[ptr[u]]
                v = hacia[e]
                if cap[e] and nivel[v] == nivel[u] + 1 and dfs(v):
                    cap[e] -= 1
                    cap[e ^ 1] += 1       # el reverso gana capacidad
                    return True
                ptr[u] += 1
            return False
        while dfs(s):
            flujo += 1
    # ---- corte: lo alcanzable desde s en el residual es un lado ----
    alcanza = [False]*(n + 1)
    alcanza[s] = True
    pila = [s]
    while pila:
        u = pila.pop()
        for e in ady[u]:
            if cap[e] and not alcanza[hacia[e]]:
                alcanza[hacia[e]] = True
                pila.append(hacia[e])
    # calles con un extremo de cada lado: estan saturadas y suman el flujo
    corte = [f"{a} {b}" for a, b in zip(A, B) if alcanza[a] != alcanza[b]]
    sys.stdout.write(str(len(corte)) + "\n" + "\n".join(corte) + ("\n" if corte else ""))

solve()
\end{lstlisting}
